\chapter{Background}
\label{ch:background}

\section{Autoencoders}

An autoencoder is a pair of functions trained together: an encoder that squeezes an input into a smaller
code, and a decoder that rebuilds the input from the code~\cite{hinton2006reducing}. Nothing tells the
encoder what the code should mean. The only pressure is that the decoder has to succeed, so the code ends up
holding whatever the decoder needs most.

That last sentence is the key to this whole report. An autoencoder keeps what helps reconstruction and is
indifferent to everything else. If two different inputs can be rebuilt almost equally well from the same
code, the encoder has no reason to tell them apart.

\section{NLA Design}

A natural language autoencoder~\cite{frasertaliente2026nla} is an autoencoder whose input is an activation
vector from a language model and whose code is a paragraph of English. Figure~\ref{fig:nla} shows the three
models involved.

\begin{figure}[t]
\centering
\begin{tikzpicture}[
  box/.style={draw, rounded corners=3pt, minimum height=1.0cm, minimum width=2.9cm, align=center, font=\small},
  data/.style={align=center, font=\small},
  arr/.style={-{Stealth[length=2.2mm]}, thick}]
\node[box] (target) {Target model\\(frozen)};
\node[data, left=0.6cm of target] (text) {passage};
\node[data, right=0.9cm of target] (h) {activation\\$h \in \mathbb{R}^{3584}$};
\node[box, right=0.9cm of h] (av) {Verbalizer};
\node[box, below=1.1cm of av] (ar) {Reconstructor};
\node[data, right=0.9cm of ar] (z) {explanation $z$\\(text)};
\node[data] (hh) at (h |- ar) {rebuilt\\activation $\hat h$};
\draw[arr] (text) -- (target);
\draw[arr] (target) -- (h);
\draw[arr] (h) -- (av);
\draw[arr] (av.east) -| (z.north);
\draw[arr] (z) -- (ar);
\draw[arr] (ar) -- (hh);
\draw[{Stealth[length=2mm]}-{Stealth[length=2mm]}, dashed] (h) -- node[left, font=\small]{$\cos(h,\hat h)$} (hh);
\end{tikzpicture}
\caption[The three models in a natural language autoencoder]{The three models in a natural language
autoencoder. The target model reads a passage; one activation $h$ is taken from its residual stream. The
verbalizer turns $h$ into text, and the reconstructor turns the text back into a vector. Training rewards a
high cosine between $h$ and $\hat h$.}
\label{fig:nla}
\end{figure}

\begin{description}
\item[Target model.] The model being explained. Here it is Qwen2.5-7B-Instruct~\cite{qwen25}, and the
activation is the residual stream after block 20 of 28, a vector of 3{,}584 numbers at one token position.
\item[Verbalizer.] A copy of the target model, fine-tuned to describe an activation. The vector is rescaled
and placed into a fixed prompt as if it were the embedding of a single token; the verbalizer then writes an
explanation, one token at a time. This is the encoder.
\item[Reconstructor.] Another fine-tuned copy, cut off after the extraction layer, with a linear layer on
top. It reads the explanation and outputs a vector. This is the decoder.
\end{description}

Both vectors are scaled to unit length before they are compared, so the reconstruction error is
$2(1-\cos(h,\hat h))$: only the direction matters. I report the cosine throughout.

The authors train in three stages. The reconstructor first learns to predict activations from text. The
verbalizer first learns to imitate explanations that a large language model, called through an API, wrote
for each position. Then the two are trained together: the verbalizer by reinforcement learning
(GRPO~\cite{shao2024deepseekmath}), with the reconstruction score as its reward, and the reconstructor by
ordinary regression on whatever the verbalizer writes. In this last stage the only reward is reconstruction.

The explanations the released verbalizer writes have a steady three-part shape, visible in the example in
Chapter~\ref{ch:intro}. The first paragraph gives the genre and topic of the text. The second quotes or
paraphrases the sentence being read. The third describes the final token and lists what is expected next. I
use this shape later, to ask whether wrong names sit in the part that describes the context or the part that
predicts the continuation.

\section{Specifics}

By a \emph{specific} I mean a named entity that appears in the passage the target model read: a particular
country, city, person, animal, food or instrument. An explanation can do three things with it.

\begin{itemize}
\item Name it. I call this a \emph{true name}.
\item Name a different entity of the same type that is not in the passage, as with the piano above. I call
this a \emph{wrong name}.
\item Name nothing of that type, and describe the passage only in general terms.
\end{itemize}

The third option is not a failure. If the activation no longer holds the entity, the faithful thing to say
is nothing. The failure this report studies is the second one, because a wrong name is indistinguishable from
a true one to anyone reading the explanation without the passage in front of them.

Specifics are the natural place to look for trouble. The gist of a passage (music, a biography, a list of
recordings) is shared by thousands of training examples and moves the activation a long way, so the
reconstructor depends on it and the verbalizer learns to report it. Which instrument it was moves the
activation much less. Chapter~\ref{ch:incentives} measures exactly how much less.
